\chapter{Finite-horizon MDPs: trajectories, occupancy measures and variances}
\label{chap:pre_mdp}

This chapter develops the finite-horizon MDP theory that the analysis of Entropic-BPI uses
beyond the model of Chapter~\ref{chap:mdp}: the iterated Markov property of trajectory laws,
occupancy measures, the unrolling of a backward recursive inequality along the trajectory
(the step of Appendix~B.1 of \cite{essakine2026tight} that turns the certificate recursion
into an expectation over the episode), a Cauchy--Schwarz inequality along trajectories, the
entropic variances and their Bellman recursion (the paper's Lemma~23), the telescoping
identity that expresses the variance of the exponentiated return as a sum of one-step
variances, and the sampling structure at the visits of a state-action pair in a run of an
arbitrary algorithm (what the concentration events of Chapter~\ref{chap:empirical} consume).
Everything is stated for an arbitrary finite-horizon MDP and policy, in library generality
(namespace \texttt{Learning.MDP}, files \texttt{Episodic.lean}, \texttt{Entropic.lean} and
\texttt{Empirical.lean} of
\texttt{Essakine2026Tight/LeanMachineLearning/ReinforcementLearning/MDP/}).

\textbf{Standing notation.} $\mathcal M = (\Ss, \As, H, (p_h)_{h \le H}, (R_h)_{h \le H})$ is
a finite-horizon MDP (Definition~\ref{def:episodic_mdp}), $\pi$ a deterministic policy
(Definition~\ref{def:policy}), $s_1 \in \Ss$ an initial state and $\beta \in \R$. For
$1 \le h \le H + 1$ and $s \in \Ss$, $\Prob^\pi_{(s,h)}$ is the law of the trajectory
$(S_i, A_i, R_i)_{i \ge h}$ of $\pi$ started at $s$ at step $h$ (Definition~\ref{def:step_law}),
$\E^\pi_{(s,h)}$ its expectation, $R^\pi_h := \sum_{i = h}^H R_i$ the return from step $h$
(Definition~\ref{def:episode_return}, $R^\pi_{H+1} = 0$), and $\Prob^\pi := \Prob^\pi_{(s_1, 1)}$,
$\E^\pi := \E^\pi_{(s_1,1)}$. For $f : \Ss \to \R$ we write
\[
  (p_h f)(s, a) := \sum_{s' \in \Ss} p_h(s' \mid s, a)\, f(s'), \qquad
  \Var_{p_h}(f)(s, a) := \sum_{s'} p_h(s' \mid s, a)\big(f(s') - (p_h f)(s,a)\big)^2
  = (p_h f^2)(s,a) - \big((p_h f)(s,a)\big)^2 ,
\]
the mean and the variance of $f$ under the transition law $p_h(\cdot \mid s, a)$ (Lean:
\texttt{vecExp}, \texttt{vecVar} in \texttt{Vec.lean}). All functions on the finite sets
$\Ss$, $\As$, $\Ss \times \As$ are bounded, hence integrable against any probability measure.

\textbf{The Markov property.} We use Lemma~\ref{lem:step_law_markov} in the following form.
For $1 \le h \le H$ and $s \in \Ss$: under $\Prob^\pi_{(s,h)}$, $S_h = s$ and $A_h = \pi_h(s)$
almost surely, $(R_h, S_{h+1}) \sim R_h(\cdot \mid s, \pi_h(s)) \otimes p_h(\cdot \mid s, \pi_h(s))$,
and conditionally on the first round $(S_h, A_h, R_h, S_{h+1})$ the shifted trajectory
$\theta := (S_i, A_i, R_i)_{i \ge h+1}$ has law $\Prob^\pi_{(S_{h+1}, h+1)}$; in integrated
form, for bounded measurable $\psi : \As \times \R \times \Ss \to \R$ and $F$ on trajectories,
\begin{equation}\label{eq:pmdp_markov_integrated}
  \E^\pi_{(s,h)}\big[\psi(A_h, R_h, S_{h+1})\, F(\theta)\big]
  = \E_{(R, s') \sim (R_h \otimes p_h)(s, \pi_h(s))}\Big[\psi(\pi_h(s), R, s')\;\E^\pi_{(s',h+1)}[F]\Big] .
\end{equation}
For $h = H + 1$ the trajectory law $\Prob^\pi_{(s, H+1)}$ is the law of the absorbing
terminal layer (all rewards $0$), and $R^\pi_{H+1} = 0$.

\textbf{What Mathlib/LML provides.} LML's \texttt{MDP.policyMeasure} and its Ionescu-Tulcea
structure (\texttt{trajMeasure}, \texttt{IsAlgEnvSeq}, \texttt{hasCondDistrib\_feedback},
\texttt{hasCondDistrib\_action}), the stationary environment \texttt{stationaryEnv} with
\texttt{feedback\_stationaryEnv}, and, for bandits, the law of the reward of the $k$-th pull
of an arm (\texttt{Online/Bandit/RewardByCountMeasure.lean}: \texttt{hasLaw\_rewardByCount},
\texttt{iIndepFun\_rewardByCount}). Mathlib: \texttt{ProbabilityTheory.variance},
\texttt{variance\_def'}, \texttt{MeasureTheory.integral\_finset\_sum},
\texttt{integral\_indicator}, \texttt{Nat.nth} for the enumeration of the visits, and the
Cauchy--Schwarz inequalities \texttt{Finset.sum\_mul\_sq\_le\_sq\_mul\_sq} and
\texttt{integral\_mul\_le\_Lp\_mul\_Lq\_of\_nonneg}.

\textbf{What is missing.} Everything below: the iterated Markov property for finite-horizon
trajectory laws, occupancy measures, the unrolling lemma, entropic variances and the
sampling-at-visits lemma (the episodic analogue of \texttt{RewardByCountMeasure.lean}, a
candidate for LML).

\section{The iterated Markov property}
\label{sec:pre_mdp_markov}

\begin{lemma}[Iterated Markov property]\label{lem:pmdp_markov_iter}
\uses{def:episodic_mdp, def:policy, def:step_law}
Let $1 \le h' \le h \le H + 1$ and $s \in \Ss$. For every bounded measurable function $\Psi$ of
the rounds $(S_i, A_i, R_i)_{h' \le i < h}$ and every bounded measurable function $F$ of the
trajectory $(S_i, A_i, R_i)_{i \ge h}$,
\[
  \E^\pi_{(s,h')}\big[\Psi\cdot F\big] = \E^\pi_{(s,h')}\big[\Psi\cdot \Phi_F(S_h)\big],
  \qquad \Phi_F(s') := \E^\pi_{(s', h)}[F] ,
\]
i.e.\ conditionally on the rounds before $h$, the trajectory from $h$ has law
$\Prob^\pi_{(S_h, h)}$. Consequently, for $h \le H$, bounded measurable
$\psi : \As \times \R \times \Ss \to \R$ and every $f : \Ss \to \R$,
\begin{align*}
  \E^\pi_{(s,h')}\big[\Psi\cdot\psi(A_h, R_h, S_{h+1})\big]
  &= \E^\pi_{(s,h')}\Big[\Psi\cdot\E_{(R, s'') \sim (R_h \otimes p_h)(S_h, \pi_h(S_h))}\big[\psi(\pi_h(S_h), R, s'')\big]\Big], \\
  \E^\pi_{(s,h')}\big[\Psi\cdot f(S_{h+1})\big]
  &= \E^\pi_{(s,h')}\big[\Psi\cdot (p_h f)(S_h, \pi_h(S_h))\big] .
\end{align*}
\end{lemma}

\begin{proof}
\uses{lem:step_law_markov}
Induction on $k := h - h' \ge 0$, for all $h'$ and $s$ simultaneously. For $k = 0$ there are
no rounds before $h = h'$, so $\Psi$ is a constant $c$; since $S_{h'} = s$ a.s.\ under
$\Prob^\pi_{(s,h')}$ (Lemma~\ref{lem:step_law_markov}), $\Phi_F(S_{h'}) = \Phi_F(s) =
\E^\pi_{(s,h')}[F]$ a.s., and both sides equal $c\,\E^\pi_{(s,h')}[F]$. For the step, let
$h = h' + k + 1 \le H + 1$ (so $h' \le H$). Write the rounds before $h$ as the first round
$(S_{h'}, A_{h'}, R_{h'}, S_{h'+1})$ followed by the rounds $h' + 1 \le i < h$ of the shifted
trajectory $\theta$, and apply~\eqref{eq:pmdp_markov_integrated} to the function
$(A_{h'}, R_{h'}, S_{h'+1}, \theta) \mapsto \Psi\cdot F$:
\[
  \E^\pi_{(s,h')}[\Psi F]
  = \E_{(R,s')}\Big[\E^\pi_{(s',h'+1)}\big[\Psi_{R,s'}\cdot F\big]\Big],
\]
where $(R, s') \sim (R_{h'} \otimes p_{h'})(s, \pi_{h'}(s))$ and $\Psi_{R,s'}$ denotes $\Psi$
with its first round frozen at $(s, \pi_{h'}(s), R, s')$, a bounded measurable function of the
rounds $h' + 1 \le i < h$. By the induction hypothesis (with $h' + 1$ in place of $h'$,
$h - (h' + 1) = k$, and the starting state $s'$),
$\E^\pi_{(s',h'+1)}[\Psi_{R,s'} F] = \E^\pi_{(s',h'+1)}[\Psi_{R,s'}\Phi_F(S_h)]$. Applying
\eqref{eq:pmdp_markov_integrated} again, now to the function
$(A_{h'}, R_{h'}, S_{h'+1}, \theta) \mapsto \Psi\cdot\Phi_F(S_h)$, gives
$\E_{(R,s')}[\E^\pi_{(s',h'+1)}[\Psi_{R,s'}\Phi_F(S_h)]] = \E^\pi_{(s,h')}[\Psi\,\Phi_F(S_h)]$,
which proves the claim. For the consequences, take $F := \psi(A_h, R_h, S_{h+1})$, a function
of the first round of the trajectory from $h$: by Lemma~\ref{lem:step_law_markov} under
$\Prob^\pi_{(s',h)}$, $\Phi_F(s') = \E_{(R,s'') \sim (R_h\otimes p_h)(s', \pi_h(s'))}[\psi(\pi_h(s'), R, s'')]$;
with $\psi(a, r, s'') := f(s'')$ this is $\sum_{s''} p_h(s'' \mid s', \pi_h(s')) f(s'') = (p_h f)(s', \pi_h(s'))$.
\end{proof}

\textbf{Lean remark.} In \texttt{stepLaw M π (s, h)} the trajectory is a sequence of rounds
indexed by $\N$ and the paper's step $i$ is the round $i - h'$ (with the $0$-indexed
translation of Chapter~\ref{chap:mdp}); ``the trajectory from $h$'' is the shift by $h - h'$
rounds, and ``the rounds before $h$'' are the first $h - h'$ rounds. State the lemma as a
\texttt{HasCondDistrib} statement (the conditional law of the shifted trajectory given the
first $h - h'$ rounds is \texttt{stepLaw M π (S\_h, h)}), or directly in the integral form
above, by induction on $h - h'$ from LML's one-step statement. The layer component of the
state at round $k$ is deterministic ($h' + k$); this is part of the description of
\texttt{layerMDP} and is what allows writing $\pi_h(S_h)$ for the action.

\begin{lemma}[Actions and rewards along a trajectory]\label{lem:pmdp_action_eq_policy}
\uses{def:episodic_mdp, def:policy, def:step_law, def:reward_fn}
Let $1 \le h' \le h \le H$ and $s \in \Ss$. Under $\Prob^\pi_{(s,h')}$: $S_{h'} = s$ and
$A_h = \pi_h(S_h)$ almost surely; and if $\mathcal M$ has the reward function $r$
(Definition~\ref{def:reward_fn}) then $R_h = r_h(S_h, A_h)$ almost surely.
\end{lemma}

\begin{proof}
\uses{lem:step_law_markov, lem:pmdp_markov_iter}
$S_{h'} = s$ a.s.\ is part of Lemma~\ref{lem:step_law_markov}. Apply
Lemma~\ref{lem:pmdp_markov_iter} with $\Psi = 1$ and $F := \indic\{A_h \ne \pi_h(S_h)\}$, a
function of the first round of the trajectory from $h$: $\Phi_F(s') = \Prob^\pi_{(s',h)}(A_h \ne \pi_h(S_h)) = 0$
for every $s'$, because $S_h = s'$ and $A_h = \pi_h(s')$ a.s.\ under $\Prob^\pi_{(s',h)}$
(Lemma~\ref{lem:step_law_markov}). Hence $\Prob^\pi_{(s,h')}(A_h \ne \pi_h(S_h)) = 0$. For the
rewards, take $F := \indic\{R_h \ne r_h(S_h, A_h)\}$: under $\Prob^\pi_{(s',h)}$, $R_h$ has law
$R_h(\cdot \mid s', \pi_h(s')) = \delta_{r_h(s', \pi_h(s'))}$, so again $\Phi_F \equiv 0$.
\end{proof}

\section{Occupancy measures}
\label{sec:pre_mdp_occupancy}

\begin{definition}[Occupancy measure]\label{def:occupancy}
\lean{Learning.MDP.Episodic.occupancy}\leanok
\uses{def:episodic_mdp, def:policy, def:step_law}
For $1 \le h \le H$, $s \in \Ss$ and $a \in \As$, the \emph{occupancy measure} of $\pi$ from
$s_1$ at step $h$ is
\[
  p^\pi_h(s, a) := \Prob^\pi(S_h = s,\ A_h = a) = \Prob^\pi_{(s_1,1)}(S_h = s,\ A_h = a) .
\]
\end{definition}

\textbf{Lean remark.} \texttt{occupancy M π s₁ h s a := (stepLaw M π (s₁, 0)).real \{ω | (ω h).obs.1 = s ∧ (ω h).action = a\}}
with the $0$-indexed step \texttt{h : Fin H} (the paper's $h$ is Lean's $h - 1$), a real number
(\texttt{Measure.real}). The dependence on $s_1$ is kept explicit.

\begin{lemma}[Occupancy measures are probability vectors]\label{lem:occupancy_sum_one}
\uses{def:occupancy}
For $1 \le h \le H$: $p^\pi_h(s,a) \ge 0$ for all $(s, a)$, $\sum_{s,a} p^\pi_h(s,a) = 1$, and
\[
  p^\pi_h(s, a) = \indic\{a = \pi_h(s)\}\,\Prob^\pi(S_h = s) .
\]
\end{lemma}

\begin{proof}
\uses{lem:pmdp_action_eq_policy}
Nonnegativity is clear. The events $\{S_h = s, A_h = a\}$, $(s, a) \in \Ss \times \As$, are
pairwise disjoint and cover $\Omega$, so their probabilities sum to $1$ (finite additivity,
Mathlib \texttt{measureReal\_biUnion\_finset} or \texttt{sum\_measureReal\_preimage\_singleton}
for the map $\omega \mapsto (S_h, A_h)$). By Lemma~\ref{lem:pmdp_action_eq_policy},
$A_h = \pi_h(S_h)$ a.s., so up to a null set $\{S_h = s, A_h = a\} = \{S_h = s, \pi_h(s) = a\}$,
which is $\{S_h = s\}$ if $a = \pi_h(s)$ and empty otherwise.
\end{proof}

\begin{lemma}[Expectations of sums along the trajectory]\label{lem:occupancy_expectation}
\uses{def:occupancy}
For functions $f_h : \Ss \times \As \to \R$, $1 \le h \le H$,
\[
  \E^\pi\Big[\sum_{h=1}^H f_h(S_h, A_h)\Big] = \sum_{h=1}^H\sum_{s,a} p^\pi_h(s,a)\, f_h(s,a) .
\]
\end{lemma}

\begin{proof}
\uses{def:occupancy}
Pointwise, $f_h(S_h, A_h) = \sum_{s,a} f_h(s,a)\,\indic\{S_h = s, A_h = a\}$ (exactly one term
of the sum is nonzero). All functions are bounded, so the integral commutes with the finite
sums (\texttt{integral\_finset\_sum}, \texttt{integral\_const\_mul}), and
$\E^\pi[\indic\{S_h = s, A_h = a\}] = p^\pi_h(s,a)$ (\texttt{integral\_indicator\_one} or
\texttt{integral\_indicator} with \texttt{Measure.real}).
\end{proof}

\begin{lemma}[Forward recursion of the occupancy measures]\label{lem:occupancy_recursion}
\uses{def:occupancy}
For $1 \le h \le H - 1$ and $(s', a') \in \Ss \times \As$,
\[
  p^\pi_{h+1}(s', a') = \indic\{a' = \pi_{h+1}(s')\}\sum_{s,a} p^\pi_h(s,a)\, p_h(s' \mid s, a) .
\]
\end{lemma}

\begin{proof}
\uses{lem:occupancy_sum_one, lem:pmdp_markov_iter}
By Lemma~\ref{lem:occupancy_sum_one} at step $h + 1$ it suffices to compute
$\Prob^\pi(S_{h+1} = s')$. Decomposing over the value of $S_h$,
\[
  \Prob^\pi(S_{h+1} = s') = \sum_{s}\E^\pi\big[\indic\{S_h = s\}\,\indic\{S_{h+1} = s'\}\big]
  = \sum_{s}\E^\pi\big[\indic\{S_h = s\}\,p_h(s' \mid S_h, \pi_h(S_h))\big]
  = \sum_{s}\Prob^\pi(S_h = s)\,p_h(s' \mid s, \pi_h(s)) ,
\]
where the second equality is Lemma~\ref{lem:pmdp_markov_iter} (with $h' = 1$,
$\Psi = \indic\{S_h = s\}$ and $f = \indic\{\cdot = s'\}$, so that $(p_h f)(s, \pi_h(s)) = p_h(s' \mid s, \pi_h(s))$).
By Lemma~\ref{lem:occupancy_sum_one}, $\Prob^\pi(S_h = s)\,p_h(s'\mid s, \pi_h(s)) = \sum_a p^\pi_h(s,a)\,p_h(s' \mid s, a)$
(the terms with $a \ne \pi_h(s)$ vanish), which gives the claim.
\end{proof}

\section{Unrolling a backward recursion along the trajectory}
\label{sec:pre_mdp_unroll}

The certificate of Entropic-BPI satisfies, on the good event, a backward inequality of the
form $g_h(s) \le e^{\beta r_h(s,a)}[u_h(s,a) + \kappa\,(p_h g_{h+1})(s,a)]$ at the action
$a = \pi_h(s)$ of the greedy policy (Lemma~\ref{lem:cert_true_recursion_pos}). The next lemma
unrolls such an inequality into an expectation along the trajectory of $\pi$; it is the
unrolling step of Appendix~B.1 of \cite{essakine2026tight} (``Sample complexity''), stated for
arbitrary functions.

\begin{lemma}[Unrolling a backward recursive inequality]\label{lem:unroll_recursion}
\uses{def:episodic_mdp, def:policy, def:step_law}
Let $\beta \in \R$, $\kappa \ge 0$, and let $r_h, u_h : \Ss \times \As \to \R$ ($1 \le h \le H$)
and $g_h : \Ss \to \R$ ($1 \le h \le H + 1$) be arbitrary functions with $g_{H+1} = 0$. Assume
that for all $1 \le h \le H$ and $s \in \Ss$, with $a := \pi_h(s)$,
\[
  g_h(s) \le e^{\beta r_h(s,a)}\Big[u_h(s,a) + \kappa\,(p_h g_{h+1})(s,a)\Big] .
\]
Then for every $1 \le h \le H + 1$,
\begin{equation}\label{eq:pmdp_unroll_ih}
  \E^\pi\Big[e^{\beta\sum_{i < h} r_i(S_i, A_i)}\,g_h(S_h)\Big]
  \le \sum_{h' = h}^{H}\kappa^{h' - h}\,\E^\pi\Big[e^{\beta\sum_{i \le h'} r_i(S_i,A_i)}\,u_{h'}(S_{h'}, A_{h'})\Big] ,
\end{equation}
and in particular ($h = 1$, using $S_1 = s_1$ a.s.)
\[
  g_1(s_1) \le \sum_{h = 1}^{H}\kappa^{h - 1}\,\E^\pi\Big[e^{\beta\sum_{i \le h} r_i(S_i,A_i)}\,u_h(S_h, A_h)\Big] .
\]
\end{lemma}

\begin{proof}
\uses{lem:pmdp_action_eq_policy, lem:pmdp_markov_iter}
Write $W_h := e^{\beta\sum_{i < h} r_i(S_i, A_i)}$ for $1 \le h \le H + 1$, a positive function
of the rounds before $h$, with $W_1 = 1$ and $W_{h+1} = W_h\,e^{\beta r_h(S_h, A_h)}$. We prove
\eqref{eq:pmdp_unroll_ih} by backward induction on $h$.

\emph{Base case $h = H + 1$.} The left-hand side is $\E^\pi[W_{H+1}\,g_{H+1}(S_{H+1})] = 0$
since $g_{H+1} = 0$, and the right-hand side is an empty sum, equal to $0$.

\emph{Induction step.} Let $1 \le h \le H$ and assume \eqref{eq:pmdp_unroll_ih} for $h + 1$.
By Lemma~\ref{lem:pmdp_action_eq_policy}, $A_h = \pi_h(S_h)$ a.s., so the hypothesis applied at
$s = S_h$ gives, almost surely,
\[
  g_h(S_h) \le e^{\beta r_h(S_h, A_h)}\Big[u_h(S_h, A_h) + \kappa\,(p_h g_{h+1})(S_h, A_h)\Big] .
\]
Multiplying by $W_h > 0$, using $W_h e^{\beta r_h(S_h,A_h)} = W_{h+1}$, and taking expectations
(all functions are bounded):
\[
  \E^\pi[W_h\,g_h(S_h)] \le \E^\pi\big[W_{h+1}\,u_h(S_h, A_h)\big]
  + \kappa\,\E^\pi\big[W_{h+1}\,(p_h g_{h+1})(S_h, A_h)\big] .
\]
The first term is the summand $h' = h$ of the right-hand side of~\eqref{eq:pmdp_unroll_ih}
(with $\kappa^0 = 1$). For the second, $W_{h+1}$ is a function of the rounds before $h + 1$
and $A_h = \pi_h(S_h)$ a.s., so Lemma~\ref{lem:pmdp_markov_iter} (with $\Psi = W_{h+1}$ and
$f = g_{h+1}$) gives $\E^\pi[W_{h+1}\,(p_h g_{h+1})(S_h, A_h)] = \E^\pi[W_{h+1}\,g_{h+1}(S_{h+1})]$,
which by the induction hypothesis is at most
$\sum_{h' = h+1}^{H}\kappa^{h' - h - 1}\E^\pi[e^{\beta\sum_{i \le h'} r_i}u_{h'}(S_{h'},A_{h'})]$.
Multiplying by $\kappa \ge 0$ turns $\kappa^{h' - h - 1}$ into $\kappa^{h' - h}$, and adding the
summand $h' = h$ gives \eqref{eq:pmdp_unroll_ih} for $h$.

For the last claim, take $h = 1$: $W_1 = 1$ and $S_1 = s_1$ a.s.\ (Lemma~\ref{lem:pmdp_action_eq_policy}),
so the left-hand side is $g_1(s_1)$.
\end{proof}

\begin{remark}[Hypotheses]\label{rem:pmdp_unroll_hypotheses}
The paper applies the unrolling with $\kappa = 1 + 13/H$, nonnegative $u_h$ and rewards in
$[0, 1]$; none of $\kappa \ge 1$, $u_h \ge 0$, $r_h \ge 0$ is needed, only $\kappa \ge 0$ (to
multiply an inequality) and $e^{\beta r} > 0$ (to keep the direction of the inequality when
multiplying by $W_h$). The functions $r_h$ are arbitrary: the lemma is about the state-action
process only, and the reward kernels of $\mathcal M$ play no role. \textbf{Lean remark.} The
induction is on the number $j = H + 1 - h$ of steps to the end (the convention of the backward
recursions of Chapter~\ref{chap:algorithm}); $\sum_{i \le h} r_i(S_i, A_i)$ is a
\texttt{Finset.sum} over \texttt{Finset.range h} of the rounds of the trajectory, and
\eqref{eq:pmdp_unroll_ih} is the statement to prove by induction (the ``in particular'' is its
instance $h = 1$).
\end{remark}

\begin{lemma}[Cauchy--Schwarz along a trajectory]\label{lem:cauchy_schwarz_traj}
Let $(\Omega, \Prob)$ be a probability space and, for $1 \le h \le H$, let
$W_h, V_h, U_h : \Omega \to [0, \infty)$ be bounded random variables. Then
\[
  \E\Big[\sum_{h = 1}^H W_h\sqrt{V_h U_h}\Big]
  \le \sqrt{\E\Big[\sum_{h=1}^H W_h^2 V_h\Big]}\;\sqrt{\E\Big[\sum_{h=1}^H U_h\Big]} .
\]
In particular this applies to $W_h = w_h((S_i, A_i)_{i \le h})$, $V_h = v_h(S_h, A_h)$,
$U_h = u_h(S_h, A_h)$ under $\Prob^\pi$ for nonnegative functions $w_h, v_h, u_h$.
\end{lemma}

\begin{proof}
Let $X_h := W_h\sqrt{V_h} \ge 0$ and $Y_h := \sqrt{U_h} \ge 0$, so that
$X_h Y_h = W_h\sqrt{V_h U_h}$ (\texttt{Real.sqrt\_mul}), $X_h^2 = W_h^2 V_h$ and $Y_h^2 = U_h$
(\texttt{Real.sq\_sqrt}). By the Cauchy--Schwarz inequality for integrals (Mathlib
\texttt{integral\_mul\_le\_Lp\_mul\_Lq\_of\_nonneg} with $p = q = 2$; the variables are bounded,
hence in $L^2$), $\E[X_h Y_h] \le \sqrt{\E X_h^2}\sqrt{\E Y_h^2}$ for each $h$. Summing over
$h$ and applying the Cauchy--Schwarz inequality for finite sums to the vectors
$(\sqrt{\E X_h^2})_h$ and $(\sqrt{\E Y_h^2})_h$ (\texttt{Finset.sum\_mul\_sq\_le\_sq\_mul\_sq},
then \texttt{Real.sqrt\_le\_sqrt} and \texttt{Real.sqrt\_mul}),
\[
  \E\Big[\sum_h X_h Y_h\Big] = \sum_h\E[X_h Y_h] \le \sum_h\sqrt{\E X_h^2}\sqrt{\E Y_h^2}
  \le \sqrt{\sum_h\E X_h^2}\,\sqrt{\sum_h\E Y_h^2}
  = \sqrt{\E\Big[\sum_h X_h^2\Big]}\,\sqrt{\E\Big[\sum_h Y_h^2\Big]} ,
\]
using the linearity of the integral for the finite sums (\texttt{integral\_finset\_sum}).
\end{proof}

\textbf{Lean remark.} The lemma is pure probability (no MDP); it is used in
Lemma~\ref{lem:ratio_bound_pos} with $W_h = e^{\beta\sum_{i \le h} r_i(S_i,A_i)}$, which is
a function of the whole prefix of the trajectory, not of $(S_h, A_h)$ alone: this is why the
statement is for general random variables and not through the occupancy measures.

\section{Entropic variances}
\label{sec:pre_mdp_entropic_variance}

Throughout this section the reward kernels of $\mathcal M$ are supported in a bounded
interval (the paper's rewards are deterministic with values in $[0, 1]$), so that every
exponentiated return $e^{\beta R^\pi_h}$ is bounded and square-integrable, and for the
recursion (Lemma~\ref{lem:entropic_variance_recursion} onwards) $\mathcal M$ has a
deterministic reward function $r$ (Definition~\ref{def:reward_fn}).

\begin{definition}[Entropic variances]\label{def:entropic_variance}
\lean{Learning.MDP.Episodic.entropicVarQ, Learning.MDP.Episodic.entropicVarV}\leanok
\uses{def:episodic_mdp, def:policy, def:step_law, def:episode_return, def:exp_value}
For $1 \le h \le H$ and $(s, a) \in \Ss \times \As$, let $\Prob^\pi_{(s,a,h)}$ be the law of
$(R, (S_i, A_i, R_i)_{i \ge h+1})$ where $(R, S') \sim (R_h \otimes p_h)(s, a)$ and, given
$(R, S')$, the trajectory from step $h + 1$ has law $\Prob^\pi_{(S', h+1)}$: the law of the
rest of the episode when the action $a$ is played at $s$ at step $h$ and $\pi$ is followed
afterwards. Under it, $R + R^\pi_{h+1}$ is the return from $(s, a)$, and by
Definition~\ref{def:exp_value} and Fubini,
$\E^\pi_{(s,a,h)}\big[e^{\beta(R + R^\pi_{h+1})}\big] = \E\big[e^{\beta R}Z^\pi_{h+1}(S')\big] = U^\pi_h(s,a)$.
The \emph{entropic variances} of $\pi$ are
\[
  \sigma Q^\pi_h(s, a) := \Var_{\Prob^\pi_{(s,a,h)}}\big(e^{\beta(R + R^\pi_{h+1})}\big)
  = \E^\pi_{(s,a,h)}\Big[\big(e^{\beta(R + R^\pi_{h+1})} - U^\pi_h(s,a)\big)^2\Big],
  \qquad
  \sigma V^\pi_h(s) := \sigma Q^\pi_h(s, \pi_h(s)),
\]
for $1 \le h \le H$, and $\sigma V^\pi_{H+1} := 0$.
\end{definition}

\textbf{Lean remark.} \texttt{entropicVarQ M β π h s a} is
\texttt{variance (fun x ↦ exp (β * (x.1 + episodeReturn x.2))) (saLaw M π h s a)} where
\texttt{saLaw M π h s a} is the composition of the step kernel $(R_h \otimes p_h)(s,a)$ with
\texttt{fun (R, s') ↦ (dirac R).prod (stepLaw M π (s', h + 1))} (``\texttt{stepLaw} from
\texttt{(x.2, h+1)} after the first transition''); \texttt{entropicVarV M β π h s :=
entropicVarQ M β π h s (π h s)} for $h \le H$ and $0$ at the terminal step. This matches the
paper's ``$\E^\pi[\,\cdot \mid S_h = s, A_h = a]$'' (Lemma~23) without conditioning on an event
of probability zero when $a \ne \pi_h(s)$.

\begin{lemma}[The entropic variance of a state is a variance under the trajectory law]\label{lem:pmdp_entropic_variance_V}
\uses{def:entropic_variance, def:exp_value}
For $1 \le h \le H$ and $s \in \Ss$,
\[
  \sigma V^\pi_h(s) = \Var_{\Prob^\pi_{(s,h)}}\big(e^{\beta R^\pi_h}\big)
  = \E^\pi_{(s,h)}\Big[\big(e^{\beta R^\pi_h} - Z^\pi_h(s)\big)^2\Big] ,
\]
and both sides are $0$ for $h = H + 1$.
\end{lemma}

\begin{proof}
\uses{lem:step_law_markov, lem:exp_bellman}
Under $\Prob^\pi_{(s,h)}$, $A_h = \pi_h(s)$ a.s.\ and the pair $(R_h, \theta)$, with $\theta$
the trajectory from $h + 1$, has law $\Prob^\pi_{(s, \pi_h(s), h)}$ by
Lemma~\ref{lem:step_law_markov} (the law of $(R_h, S_{h+1})$ is $(R_h \otimes p_h)(s, \pi_h(s))$
and, given the first round, $\theta \sim \Prob^\pi_{(S_{h+1}, h+1)}$), and
$R^\pi_h = R_h + R^\pi_{h+1}(\theta)$. So $e^{\beta R^\pi_h}$ under $\Prob^\pi_{(s,h)}$ and
$e^{\beta(R + R^\pi_{h+1})}$ under $\Prob^\pi_{(s, \pi_h(s), h)}$ have the same law, hence the same
variance (\texttt{variance} depends only on the law: \texttt{variance\_map} /
\texttt{IdentDistrib.variance\_eq}). The mean of $e^{\beta R^\pi_h}$ is $Z^\pi_h(s)$
(Definition~\ref{def:exp_value}), which equals $U^\pi_h(s, \pi_h(s))$ by
Lemma~\ref{lem:exp_bellman}, and for a square-integrable variable the variance is the mean
square deviation from the mean (\texttt{variance\_def'} or the definition of
\texttt{variance} through \texttt{evariance}). For $h = H + 1$, $R^\pi_{H+1} = 0$, the variable
is the constant $1$ and $Z^\pi_{H+1} = 1$, so both sides are $0 = \sigma V^\pi_{H+1}$.
\end{proof}

\begin{lemma}[Entropic variance recursion]\label{lem:entropic_variance_recursion}
\uses{def:entropic_variance, def:exp_value, def:reward_fn}
Assume $\mathcal M$ has the reward function $r$. For $1 \le h \le H$ and $(s,a) \in \Ss\times\As$,
\[
  \sigma Q^\pi_h(s, a) = e^{2\beta r_h(s,a)}\,\Var_{p_h}\big(Z^\pi_{h+1}\big)(s,a)
  + e^{2\beta r_h(s,a)}\,\big(p_h\,\sigma V^\pi_{h+1}\big)(s,a) .
\]
\end{lemma}

\begin{proof}
\uses{lem:pmdp_entropic_variance_V, def:exp_value}
Write $r := r_h(s,a)$, $Z := Z^\pi_{h+1}$ and $m := (p_h Z)(s,a)$. Under $\Prob^\pi_{(s,a,h)}$ the
reward kernel is $\delta_{r}$, so $R = r$ a.s.\ and $e^{\beta(R + R^\pi_{h+1})} = e^{\beta r}X$ with
$X := e^{\beta R^\pi_{h+1}}$; hence $\sigma Q^\pi_h(s,a) = e^{2\beta r}\Var(X)$
(\texttt{variance\_mul\_const}). The law of $(S', \theta)$ under $\Prob^\pi_{(s,a,h)}$ is the
mixture $\sum_{s'} p_h(s' \mid s,a)\,\delta_{s'} \otimes \Prob^\pi_{(s',h+1)}$ and $X$ is a
function of $\theta$, so for every bounded measurable $\phi$,
$\E^\pi_{(s,a,h)}[\phi(X)] = \sum_{s'} p_h(s' \mid s,a)\,\E^\pi_{(s',h+1)}[\phi(X)]$. With
$\phi = \mathrm{id}$ this gives $\E[X] = \sum_{s'} p_h(s'\mid s,a) Z(s') = m$
(Definition~\ref{def:exp_value}). We now make the law of total variance (conditioning on
$S' = s'$) explicit. For each $s'$, since $\E^\pi_{(s',h+1)}[X] = Z(s')$,
\[
  \E^\pi_{(s',h+1)}\big[(X - m)^2\big]
  = \E^\pi_{(s',h+1)}\big[(X - Z(s'))^2\big] + 2\,(Z(s') - m)\,\E^\pi_{(s',h+1)}\big[X - Z(s')\big] + (Z(s') - m)^2
  = \sigma V^\pi_{h+1}(s') + (Z(s') - m)^2 ,
\]
by Lemma~\ref{lem:pmdp_entropic_variance_V} at step $h + 1$ (for $h = H$: $X = 1 = Z(s')$ and
$\sigma V^\pi_{H+1} = 0$, consistent). Therefore
\[
  \Var(X) = \E^\pi_{(s,a,h)}\big[(X - m)^2\big]
  = \sum_{s'} p_h(s' \mid s,a)\Big(\sigma V^\pi_{h+1}(s') + (Z(s') - m)^2\Big)
  = \big(p_h\,\sigma V^\pi_{h+1}\big)(s,a) + \Var_{p_h}(Z)(s,a) ,
\]
and multiplying by $e^{2\beta r}$ gives the claim. This is the paper's Lemma~23.
\end{proof}

\begin{lemma}[Telescoping of the entropic variances]\label{lem:entropic_variance_sum}
\uses{def:entropic_variance, def:exp_value, def:reward_fn}
Assume $\mathcal M$ has the reward function $r$. Then
\[
  \sum_{h=1}^{H}\E^\pi\Big[e^{2\beta\sum_{i \le h} r_i(S_i, A_i)}\,\Var_{p_h}\big(Z^\pi_{h+1}\big)(S_h, A_h)\Big]
  = \sigma V^\pi_1(s_1) = \Var_{\Prob^\pi}\big(e^{\beta R^\pi_1}\big) .
\]
\end{lemma}

\begin{proof}
\uses{lem:entropic_variance_recursion, lem:pmdp_entropic_variance_V, lem:pmdp_action_eq_policy, lem:pmdp_markov_iter}
Let $W_h := e^{2\beta\sum_{i < h} r_i(S_i, A_i)}$ for $1 \le h \le H + 1$ ($W_1 = 1$,
$W_{h+1} = W_h e^{2\beta r_h(S_h, A_h)}$), a bounded function of the rounds before $h$. Fix
$1 \le h \le H$. By Lemma~\ref{lem:pmdp_action_eq_policy}, $A_h = \pi_h(S_h)$ a.s., so
$\sigma V^\pi_h(S_h) = \sigma Q^\pi_h(S_h, A_h)$ a.s., and Lemma~\ref{lem:entropic_variance_recursion}
at $(S_h, A_h)$ gives, after multiplication by $W_h$ and integration,
\[
  \E^\pi\big[W_h\,\sigma V^\pi_h(S_h)\big]
  = \E^\pi\big[W_{h+1}\Var_{p_h}(Z^\pi_{h+1})(S_h, A_h)\big]
  + \E^\pi\big[W_{h+1}\,(p_h\,\sigma V^\pi_{h+1})(S_h, A_h)\big] .
\]
By Lemma~\ref{lem:pmdp_markov_iter} ($\Psi = W_{h+1}$, a function of the rounds before $h + 1$,
and $f = \sigma V^\pi_{h+1}$, with $A_h = \pi_h(S_h)$ a.s.), the last term equals
$\E^\pi[W_{h+1}\,\sigma V^\pi_{h+1}(S_{h+1})]$. Hence, with $D_h := \E^\pi[W_h\,\sigma V^\pi_h(S_h)]$
for $1 \le h \le H + 1$,
\[
  \E^\pi\big[W_{h+1}\Var_{p_h}(Z^\pi_{h+1})(S_h, A_h)\big] = D_h - D_{h+1} ,
\]
and $W_{h+1} = e^{2\beta\sum_{i \le h} r_i(S_i,A_i)}$. Summing over $h = 1, \dots, H$ telescopes
(\texttt{Finset.sum\_range\_sub'}) to $D_1 - D_{H+1}$. Now $D_{H+1} = 0$ because
$\sigma V^\pi_{H+1} = 0$, and $D_1 = \sigma V^\pi_1(s_1)$ because $W_1 = 1$ and $S_1 = s_1$ a.s.
Finally $\sigma V^\pi_1(s_1) = \Var_{\Prob^\pi}(e^{\beta R^\pi_1})$ by
Lemma~\ref{lem:pmdp_entropic_variance_V}.
\end{proof}

\begin{lemma}[Normalized variance of the exponentiated return]\label{lem:variance_ratio_le}
\uses{def:exp_value, def:episode_return, def:gmax}
Let $G \ge 0$ be such that $R^\pi_1 \in [0, G]$ almost surely under $\Prob^\pi$. Then
$Z^\pi_1(s_1) > 0$ and
\[
  \frac{\Var_{\Prob^\pi}\big(e^{\beta R^\pi_1}\big)}{Z^\pi_1(s_1)^2}
  \le \frac{(e^{\abs\beta G} - 1)^2}{4\,e^{\abs\beta G}}
  \le \frac{(e^{\abs\beta G} - 1)^2}{e^{\abs\beta G}} .
\]
In particular, if $\mathcal M$ has a reward function with values in $[0, 1]$, this holds with
$G = \Gmax(\mathcal M)$.
\end{lemma}

\begin{proof}
\uses{lem:normalized_variance_bound, lem:gmax_le_horizon, lem:pmdp_action_eq_policy}
Apply Lemma~\ref{lem:normalized_variance_bound} on $(\Omega, \Prob^\pi)$ with $f := R^\pi_1$ and
$R := G$: the mean of $Y = e^{\beta R^\pi_1}$ is $Z^\pi_1(s_1)$ (Definition~\ref{def:exp_value}),
which is positive, and the first inequality follows; the second holds because
$4e^{\abs\beta G} \ge e^{\abs\beta G} > 0$. For the last claim: by
Lemma~\ref{lem:gmax_le_horizon}, $\Gmax(\mathcal M) \ge 0$ and $R^\pi_1 \le \Gmax(\mathcal M)$
a.s.\ under $\Prob^\pi$; and $R^\pi_1 = \sum_{i \le H} r_i(S_i, A_i) \ge 0$ a.s.\ by
Lemma~\ref{lem:pmdp_action_eq_policy} since $r \ge 0$.
\end{proof}

\textbf{Lean remark.} This is where the paper's Lemma~24 enters (Lemma~\ref{lem:ratio_bound_pos}).
The bound is used with the paper's constant $V_G = (e^{\abs\beta G} - 1)^2/e^{\abs\beta G}$
(Definition~\ref{def:upper_bound_const}); the factor $4$ is slack.

\section{Sampling at the visits of a state-action pair}
\label{sec:pre_mdp_visits}

The concentration events of Chapter~\ref{chap:empirical} are about the transitions observed
at the visits of a fixed pair $(s, a)$ at a fixed step $h$, across the episodes of a run of an
arbitrary algorithm. The visits are chosen by the algorithm (a policy is played, the pair may
or may not be visited), so the observed next states are not literally an i.i.d.\ sample;
what is true is that each observed next state is drawn from $p_h(\cdot \mid s, a)$
independently of everything before it. We state this in the three forms that the events use.

\textbf{Setting and notation.} $(X, Y)$ is a run of an arbitrary algorithm in the episode
environment of $\mathcal M$ from $s_1$ (Definition~\ref{def:episode_env}), on $(\Omega, \Prob)$:
$X_t = \pi^{t+1}$ is the policy of episode $t + 1$ and $Y_t = (s^{t+1}_1, \dots, s^{t+1}_{H+1})$
its state sequence; $a^{t+1}_h := \pi^{t+1}_h(s^{t+1}_h)$; $\mathrm{hist}_t := (X_i, Y_i)_{i < t}$
is the history of the first $t$ episodes, and $n^t_h(s,a)$, $n^t_h(s,a,s')$, $\phat^t_h$ are the
counts and empirical transitions of Definitions~\ref{def:counts} and~\ref{def:emp_trans}. Fix
$1 \le h \le H$ and $(s, a) \in \Ss \times \As$. For $i \ge 1$ let
$w_i := \indic\{(s^i_h, a^i_h) = (s, a)\}$ be the indicator that episode $i$ visits $(s,a)$ at
step $h$, so that $n^t_h(s,a) = \sum_{i \le t} w_i$. The \emph{$k$-th visit time} ($k \ge 1$) is
\[
  T_k := \min\Big\{i \ge 1 : \sum_{j \le i} w_j = k\Big\} \in \{1, 2, \dots\} \cup \{\infty\},
\]
the index of the $k$-th episode with $w_i = 1$ ($T_k = \infty$ if there are fewer than $k$ such
episodes), and $W_k := s^{T_k}_{h+1}$ on $\{T_k < \infty\}$ (and $W_k := s$, say, elsewhere) is
the \emph{$k$-th observed next state}. Then $T_1 < T_2 < \cdots$ on the event where they are
finite, $\{T_k = i\}$ is determined by $\mathrm{hist}_{i-1}$, $X_{i-1}$ and $s^i_h$
($T_k = i$ iff $\sum_{j < i} w_j = k - 1$ and $w_i = 1$), and $n^t_h(s,a) \ge k$ iff $T_k \le t$.
For $t \ge 0$ let $\mathcal G_t$ be the $\sigma$-algebra generated by
$(\mathrm{hist}_t, X_t, s^{t+1}_1, \dots, s^{t+1}_h)$: the information available just before the
transition of episode $t + 1$ at step $h$; $w_{t+1}$ is $\mathcal G_t$-measurable.

\begin{lemma}[Counts and episodes]\label{lem:count_le_episodes}
\uses{def:counts}
For every $t \ge 0$ and $h \le H$: $0 \le n^t_h(s,a) \le t$ for every $(s,a)$, and
$\sum_{s,a} n^t_h(s,a) = t$.
\end{lemma}

\begin{proof}
\uses{def:counts}
$n^t_h(s,a)$ is a sum of $t$ indicators, hence between $0$ and $t$
(\texttt{Finset.sum\_le\_card\_nsmul}). Exchanging the sums,
$\sum_{s,a} n^t_h(s,a) = \sum_{i \le t}\sum_{s,a}\indic\{(s^i_h, a^i_h) = (s,a)\} = \sum_{i \le t} 1 = t$,
since for each episode $i$ exactly one pair equals $(s^i_h, a^i_h)$ (\texttt{Finset.sum\_comm},
\texttt{Finset.sum\_ite\_eq}).
\end{proof}

\begin{lemma}[The empirical transitions are the empirical distribution of the observed next states]\label{lem:emp_trans_reindex}
\uses{def:counts, def:emp_trans}
For every $t \ge 0$, let $n := n^t_h(s,a)$. Then $\{i \le t : w_i = 1\} = \{T_1, \dots, T_n\}$
with $T_1 < \dots < T_n \le t < T_{n+1}$, and for every $s' \in \Ss$,
\[
  n^t_h(s, a, s') = \sum_{k = 1}^{n}\indic\{W_k = s'\}, \qquad\text{hence, if } n \ge 1,\qquad
  \phat^t_h(s' \mid s, a) = \frac1n\sum_{k=1}^n\indic\{W_k = s'\} .
\]
In particular $\phat^t_h(\cdot \mid s,a)$ depends on $t$ only through $n^t_h(s,a)$: it is the
empirical distribution of $W_1, \dots, W_{n}$.
\end{lemma}

\begin{proof}
\uses{def:counts, def:emp_trans}
By definition, $T_k$ is the $k$-th element (in increasing order) of the set
$I := \{i \ge 1 : w_i = 1\}$, and $I \cap \{1, \dots, t\}$ has exactly $n = \sum_{i \le t} w_i$
elements; so $I \cap \{1,\dots,t\} = \{T_1, \dots, T_n\}$, $T_n \le t$ and $T_{n+1} > t$. Then
\[
  n^t_h(s,a,s') = \sum_{i \le t} w_i\,\indic\{s^i_{h+1} = s'\}
  = \sum_{i \in I \cap \{1,\dots,t\}}\indic\{s^i_{h+1} = s'\}
  = \sum_{k=1}^n\indic\{s^{T_k}_{h+1} = s'\} = \sum_{k=1}^n\indic\{W_k = s'\}
\]
(reindexing the sum by the bijection $k \mapsto T_k$, \texttt{Finset.sum\_bij} or
\texttt{Finset.sum\_map}), and $\phat^t_h(s'\mid s,a) = n^t_h(s,a,s')/n$ for $n \ge 1$
(Definition~\ref{def:emp_trans}).
\end{proof}

\textbf{Lean remark.} With the episodes $0$-indexed, $T_k - 1 = $ \texttt{Nat.nth (fun i ↦ w i = 1) (k - 1)}
(Mathlib \texttt{Nat.nth}, with \texttt{Nat.count}, \texttt{Nat.nth\_count},
\texttt{Nat.nth\_mem\_of\_lt\_card}, \texttt{Nat.count\_nth\_of\_infinite}/\texttt{count\_nth});
$\{T_k < \infty\}$ is ``$k - 1 <$ the number of visits'' or ``the set of visits is infinite''.
Alternatively, avoid $\infty$ by working with the finite enumeration
\texttt{Finset.orderEmbOfFin} of $\{i < t : w_i = 1\}$ for each $t$; the statement (i) below
about $\{T_n < \infty\}$ is then a statement about $\bigcup_t\{n^t_h(s,a) = n\}$.

\begin{lemma}[Sampling at the visits]\label{lem:visits_iid}
\uses{def:episode_env, def:counts, def:emp_trans, lem:emp_trans_reindex}
In the setting above, with $p := p_h(\cdot \mid s, a)$:
\begin{enumerate}
  \item[(i)] (\emph{Dominated law of the observed next states.}) For every $n \ge 1$ and
        $(w_1, \dots, w_n) \in \Ss^n$,
        \[
          \Prob\big(T_n < \infty,\ W_1 = w_1, \dots, W_n = w_n\big)
          = \Prob\big(T_n < \infty,\ W_1 = w_1, \dots, W_{n-1} = w_{n-1}\big)\,p(w_n)
          \le \prod_{k=1}^n p(w_k) .
        \]
        Consequently, for every $B \subseteq \Ss^n$,
        $\Prob\big(T_n < \infty,\ (W_1, \dots, W_n) \in B\big) \le p^{\otimes n}(B)$, and since
        $\{T_n < \infty\} = \{\exists t \ge 0 : n^t_h(s,a) = n\}$, for every $B$,
        \[
          \Prob\big(\exists t \ge 0 :\ n^t_h(s,a) = n \text{ and } (W_1, \dots, W_n) \in B\big)
          \le p^{\otimes n}(B) ,
        \]
        where $p^{\otimes n}$ is the law of $n$ i.i.d.\ draws from $p$. In words: on the event
        that the $n$-th visit happens, the first $n$ observed next states are dominated by an
        i.i.d.\ sample of $p$; by Lemma~\ref{lem:emp_trans_reindex}, this applies to every
        event about $\phat^t_h(\cdot \mid s,a)$ with $n^t_h(s,a) = n$, uniformly in $t$.
  \item[(ii)] (\emph{Martingale differences along the episodes.}) For every $t \ge 0$, every
        bounded $\mathcal G_t$-measurable $\Psi$ and every $f : \Ss \to \R$,
        \[
          \E\big[\Psi\, f(s^{t+1}_{h+1})\big] = \E\big[\Psi\,(p_h f)(s^{t+1}_h, a^{t+1}_h)\big] ,
          \qquad\text{in particular}\qquad
          \E\big[\Psi\, w_{t+1}\, f(s^{t+1}_{h+1})\big] = \E\big[\Psi\, w_{t+1}\big]\,(pf) ;
        \]
        i.e.\ $\E[f(s^{t+1}_{h+1}) \mid \mathcal G_t] = (p_h f)(s^{t+1}_h, a^{t+1}_h)$ a.s., and
        $w_{t+1}\big(f(s^{t+1}_{h+1}) - pf\big)$ is a martingale difference with respect to
        $(\mathcal G_t)_t$, with $w_{t+1}$ predictable.
  \item[(iii)] (\emph{Martingale differences along the visits.}) Let $k \ge 1$ and let $E$ be an
        event such that $E \cap \{T_k = t + 1\}$ is $\mathcal G_t$-measurable for every $t \ge 0$
        (an event of the information available at the $k$-th visit; for instance
        $E = \{W_1 = w_1, \dots, W_{k-1} = w_{k-1}\}$ or $E = \Omega$). Then for every $f : \Ss \to \R$,
        \[
          \E\big[\indic_{E}\,\indic\{T_k < \infty\}\, f(W_k)\big]
          = \Prob\big(E \cap \{T_k < \infty\}\big)\,(pf) .
        \]
\end{enumerate}
\end{lemma}

\begin{proof}
\uses{lem:episode_env_law, lem:pmdp_markov_iter, lem:emp_trans_reindex}
\emph{(ii).} Lemma~\ref{lem:episode_env_law} states that, conditionally on
$\mathrm{hist}_t$, $X_t$ and $s^{t+1}_1, \dots, s^{t+1}_h$, the next state $s^{t+1}_{h+1}$ has law
$p_h(\cdot \mid s^{t+1}_h, a^{t+1}_h)$: $\E[f(s^{t+1}_{h+1}) \mid \mathcal G_t] = (p_h f)(s^{t+1}_h, a^{t+1}_h)$
a.s. Multiplying by the bounded $\mathcal G_t$-measurable $\Psi$ and integrating gives the
first identity (this is also obtained directly: by Lemma~\ref{lem:episode_env_law} the
conditional law of $Y_t$ given $(\mathrm{hist}_t, X_t)$ is the law of the state sequence under
$\Prob^{X_t}_{(s_1,1)}$, and under that law Lemma~\ref{lem:pmdp_markov_iter} with $h' = 1$,
$\Psi$ a function of the rounds before $h + 1$ and $f$ gives
$\E^{\pi}[\Psi f(S_{h+1})] = \E^\pi[\Psi\,(p_h f)(S_h, \pi_h(S_h))]$; then integrate over
$(\mathrm{hist}_t, X_t)$). For the particular case, $\Psi w_{t+1}$ is bounded and
$\mathcal G_t$-measurable, and on $\{w_{t+1} = 1\}$, $(p_h f)(s^{t+1}_h, a^{t+1}_h) = pf$.

\emph{(iii).} For $t \ge 0$ let $E_t := E \cap \{T_k = t + 1\}$, which is $\mathcal G_t$-measurable
by assumption. On $E_t$, $w_{t+1} = 1$, so $(s^{t+1}_h, a^{t+1}_h) = (s, a)$ and
$W_k = s^{t+1}_{h+1}$. By (ii) with $\Psi = \indic_{E_t}$,
\[
  \E\big[\indic_{E_t}\, f(W_k)\big] = \E\big[\indic_{E_t}\, f(s^{t+1}_{h+1})\big]
  = \E\big[\indic_{E_t}\,(p_h f)(s^{t+1}_h, a^{t+1}_h)\big] = \Prob(E_t)\,(pf) .
\]
The events $E_t$, $t \ge 0$, are pairwise disjoint with union $E \cap \{T_k < \infty\}$; summing
over $t$ (countable additivity of $\Prob$ and of the integral: \texttt{integral\_iUnion} or
\texttt{tsum} of the indicators, all terms being bounded by $\sup\abs f\,\Prob(E_t)$ with
$\sum_t\Prob(E_t) \le 1$) gives the claim.

\emph{(i).} Fix $n \ge 1$ and $w_1, \dots, w_n$. The event $E := \{W_1 = w_1, \dots, W_{n-1} = w_{n-1}\}$
($E = \Omega$ if $n = 1$) satisfies the hypothesis of (iii) with $k = n$: on $\{T_n = t + 1\}$
the visit times $T_1 < \dots < T_{n-1}$ are at most $t$, so $W_1, \dots, W_{n-1}$ are functions
of $\mathrm{hist}_t$, and $\{T_n = t+1\}$ is $\mathcal G_t$-measurable. Applying (iii) with
$f = \indic\{\cdot = w_n\}$ gives the equality
$\Prob(T_n < \infty, W_1 = w_1, \dots, W_n = w_n) = \Prob(T_n < \infty, W_1 = w_1, \dots, W_{n-1} = w_{n-1})\,p(w_n)$.
For the inequality, induct on $n$: for $n = 1$ the right-hand side is $\Prob(T_1 < \infty)p(w_1) \le p(w_1)$;
for $n \ge 2$, $\{T_n < \infty\} \subseteq \{T_{n-1} < \infty\}$, so the first factor is at most
$\Prob(T_{n-1} < \infty, W_1 = w_1, \dots, W_{n-1} = w_{n-1}) \le \prod_{k < n} p(w_k)$ by the
induction hypothesis. For $B \subseteq \Ss^n$ (a finite set), sum over $(w_1,\dots,w_n) \in B$:
$\Prob(T_n < \infty, (W_1,\dots,W_n) \in B) = \sum_{w \in B}\Prob(T_n < \infty, W_{\le n} = w) \le \sum_{w \in B}\prod_k p(w_k) = p^{\otimes n}(B)$.
Finally $t \mapsto n^t_h(s,a)$ is nondecreasing, starts at $0$ and increases by at most $1$ per
episode, so it takes the value $n$ at some $t$ iff it reaches $n$, iff $T_n < \infty$
(indeed $n^{T_n}_h(s,a) = n$ by Lemma~\ref{lem:emp_trans_reindex}).
\end{proof}

\begin{remark}[Use and Lean encoding]\label{rem:pmdp_visits_iid}
Item (i) is what the KL event needs (Lemma~\ref{lem:event_kl_prob}): for the event
$\{\KL(\phat^t_h(\cdot\mid s,a)\,\|\,p) > x\}$ on $\{n^t_h(s,a) = n\}$, which by
Lemma~\ref{lem:emp_trans_reindex} is $\{(W_1,\dots,W_n) \in B\}$ with
$B = \{w \in \Ss^n : \KL(\mathrm{emp}(w)\|p) > x\}$, (i) bounds its probability, uniformly in
$t$, by the i.i.d.\ probability that Sanov's bound (Theorem~\ref{thm:sanov_hp}) controls.
Items (ii) and (iii) are what the Bernstein and count events need
(Lemmas~\ref{lem:event_bern_prob} and~\ref{lem:event_cnt_prob}): (ii) with the filtration
$(\mathcal G_t)$ of the episodes and the predictable weights $w_{t+1} \in \{0, 1\}$, and (iii)
when the time index is the number of visits. \textbf{Lean remark.} This is the episodic
analogue of LML's \texttt{Online/Bandit/RewardByCountMeasure.lean}
(\texttt{stepsUntil}, \texttt{rewardByCount}, \texttt{hasLaw\_rewardByCount},
\texttt{iIndepFun\_rewardByCount}): the ``arm'' is the triple $(h, s, a)$ (episode $i$ ``pulls''
it iff $w_i = 1$) and the ``reward'' is the next state $s^i_{h+1}$. The bandit file obtains the
exact i.i.d.\ law of the rewards by count by extending the probability space with an
independent i.i.d.\ stream (\texttt{P.prod (streamMeasure ν)}) that fills in the pulls that never
happen; the dominated form (i), stated on $\{T_n < \infty\}$ without any extension, is all the
union bound needs and is the statement to contribute to LML (\texttt{Empirical.lean}). Item
(ii) is \texttt{lem:episode\_env\_law} in integral form (\texttt{HasCondDistrib} of
$s^{t+1}_{h+1}$ given $\mathcal G_t$, from LML's \texttt{IsAlgEnvSeq.hasCondDistrib\_feedback}
and \texttt{feedback\_stationaryEnv}); the $\sigma$-algebra $\mathcal G_t$ is the comap of the
map $\omega \mapsto (\mathrm{hist}_t, X_t, (s^{t+1}_i)_{i \le h})$.
\end{remark}
